\documentclass[english,course]{lecturenotes}
% `lecturenotes` (notes/lecturenotes.cls, v3.2) requires LuaLaTeX --
% compile with `latexmk -lualatex main.tex` (see notes/README.md).

% ---------------------------------------------------------------------------
% Shared preamble for the lecture notes
% ---------------------------------------------------------------------------
%
% NOTE: main.tex uses the `lecturenotes` document class (notes/lecturenotes.cls),
% which already loads hyperref, amsmath (intlimits), amssymb, amsthm, tikz,
% mathtools, babel, etc., and already defines the theorem-like environments
% theorem/lemma/proposition/corollary/definition/example/remark (plus
% exercise/claim/conjecture/fact/problem/notation). Do not reload or redefine
% those here. The class requires LuaLaTeX and does not support inputenc/fontenc.

\usepackage{graphicx}
\usepackage{cleveref}

% Bibliography: biblatex + biber. Needed because classic BibTeX resolves .bib paths
% relative to the compiler's working directory (unreliable on Overleaf), while biblatex
% resolves \addbibresource relative to the main .tex file itself — which is what lets us
% keep bibliography.bib at the repo root and reference it as ../bibliography.bib from here.
\usepackage[backend=biber, style=numeric, sorting=none]{biblatex}
\addbibresource{../bibliography.bib}

% The class sets a narrow text block (right=2in, to leave room for \margintext),
% which leaves TeX little slack on lines carrying long unbreakable tokens --
% arXiv identifiers and DOIs in the bibliography, mostly. A little emergency
% stretch avoids overfull lines there without loosening normal body text.
\setlength{\emergencystretch}{2em}

\usetikzlibrary{shapes,positioning,arrows.meta}

% --- Cross-references to the four course parts ---
% The parts are \section's in a document whose section 1 is "General
% Information", so \ref{sec:partN} would print the section number (2,3,4,5)
% rather than the part number. \partref{n} prints the Roman numeral and still
% hyperlinks: "Part~\partref{1}", "Parts~\partref{1} and~\partref{2}".
\newcommand{\partref}[1]{%
  \ifcase#1\or
    \hyperref[sec:part1]{I}\or
    \hyperref[sec:part2]{II}\or
    \hyperref[sec:part3]{III}\or
    \hyperref[sec:part4]{IV}%
  \fi}

% --- Tensor-network notation macros (placeholders — adapt to your convention) ---
\newcommand{\tensor}[1]{\mathsf{#1}}
\newcommand{\rank}{\chi}
\newcommand{\bond}{r}

\title{Tensor network algorithms for computational fluid dynamics}
\shorttitle{TNs for CFD}

\author{Raghavendra D. Peddinti\affmark{1,*}, Stefano Pisoni\affmark{1,2}}
\affiliations{\affmark{1}\, Quantum Research Center, Technology Innovation Institute, Abu Dhabi, UAE\\\affmark{2}\, Institute for Quantum-Inspired and Quantum Optimization, Hamburg University of Technology, Germany\\ \affmark{*}raghavendra.peddinti@tii.ae
}


\conference{QTADE School}
\place{Bilbao, Spain}
\date{28}{09}{2026}
\dateend{02}{10}{2026}


% The class prints the title page and (for the `course` type) the table of
% contents automatically at \begin{document} -- no \maketitle/\tableofcontents needed.

\begin{document}

\clearpage
\section{General Information}
\label{sec:general-information}

\paragraph{What this course is about.}
Simulating a fluid means storing a field on a grid and pushing it forward in
time. Both halves of that sentence scale badly: the storage grows exponentially
with resolution and dimension, and each timestep touches all of it. These notes
develop one response to that problem --- represent the field as a
\textbf{tensor network}, and do the entire computation without ever writing the
field down in full.

\paragraph{The single idea.}
A field sampled on $2^n$ points is a vector in a space of dimension $2^n$, but a
\emph{physical} field is not an arbitrary point in that space. It is smooth, it
is multiscale, and its well-separated scales are weakly coupled. Such a vector
lies close to a low-dimensional manifold, described by $O(n\chi^2)$ numbers
rather than $2^n$. Every part of this course either exploits that fact, tests
it, or maps out where it stops being true.
\margintext{Keep the bet and its price together: exponential memory savings are
bought with a truncation error that must be monitored at every step. The
truncation never goes away.}

\paragraph{How the four parts fit together.}
\begin{itemize}
  \item \textbf{Part~\partref{1}} builds the representation: diagrammatic
    notation, matrix product states, the quantics encoding, and a first PDE
    solver.
  \item \textbf{Part~\partref{2}} adds the linear algebra --- solving
    $Ax = b$ on the manifold, and constructing a tensor train from a function
    you can only sample.
  \item \textbf{Part~\partref{3}} spends both on the Navier--Stokes
    equations: a full incompressible solver, and the empirical question of how
    far the ranks stretch.
  \item \textbf{Part~\partref{4}} inverts the problem: given data rather than
    an equation, recover the dynamics.
\end{itemize}

Parts~\partref{1} and~\partref{2} are prerequisites for the rest.
Parts~\partref{3} and~\partref{4} are independent of one another and can
be read in either order.

\paragraph{Prerequisites.}
Linear algebra --- in particular the singular value decomposition and its
optimality properties --- and a first course in numerical methods for PDEs
(finite differences, explicit and implicit time stepping, stability). No
background in quantum mechanics is assumed anywhere; where a quantum-many-body
term is unavoidable, its numerical-analysis synonym is given alongside it in
\cref{subsec:literature-survey}.

\paragraph{Notation and conventions.}
\begin{itemize}
  \item $n$ is the number of tensors in a chain, so the grid has $N = 2^n$
    points per encoded coordinate.
  \item $\chi$ is the bond dimension, written \textsc{tt-rank} in the
    numerical-analysis literature.
  \item Physical legs point down in diagrams; virtual legs run horizontally.
  \item \enquote{MPS} and \enquote{tensor train} are used interchangeably, as are
    \enquote{bond dimension} and \enquote{rank}; both vocabularies appear in the
    papers these notes draw on.
\end{itemize}

\paragraph{Margin notes.}
\margintext{Like this one.}
Notes set in the margin carry historical context, dates and attributions, and
warnings that would interrupt the argument if placed in the main text. They are
commentary --- nothing in the main development depends on them.

\paragraph{Materials.}
The companion slides, the notebooks and the shared bibliography all live in the
course repository. Every quantitative claim in these notes is measured there,
on a laptop, and you are encouraged to disagree with any of them by rerunning
the cell.

Each part has two notebooks: a \emph{walkthrough} that follows the lecture cell
by cell, and an \emph{exercise} set of roughly thirty minutes with solutions at
the bottom. Parts~\partref{1} and~\partref{2} build everything from scratch in
about five hundred lines of NumPy (\texttt{qtade\_tn.py}), because the mechanics
are the point; Parts~\partref{3} and~\partref{4} move to \texttt{quimb} for the
representation and compression, and keep the Part~\partref{2} solver, because no
tensor-network library ships a linear solver for $Ax=b$. That seam is itself
worth noticing.

\section{Part I: Tensor Network Representations}
\label{sec:part1}
\lecture{28}{09}{2026}

% Mirrors slides.qmd "Part I" outline. One subsection per outline item.

Everything in this course rests on a single bet. A field sampled on a fine grid
is an enormous vector, but it is not an \emph{arbitrary} enormous vector: it is
smooth, it is multiscale, and the coupling between well-separated scales is
weak. If that is true, the vector lies on a manifold of far lower dimension than
the space it sits in, and we should compute on the manifold rather than in the
ambient space.

Part~\partref{1} builds the machinery needed to state that bet precisely: a
diagrammatic language (\cref{subsec:penrose-notation}), a short account of where
the ideas came from (\cref{subsec:literature-survey}), the scaling wall they are
meant to break (\cref{subsec:numerical-methods-intro}), the ansatz itself
(\cref{subsec:mps}), the binary re-indexing that makes it multiscale
(\cref{subsec:qtt-functions-operators}), and --- by the end --- a working PDE
solver whose cost is logarithmic in the number of grid points
(\cref{subsec:heat-explicit-euler}).

\subsection{Penrose Diagram Notation}
\label{subsec:penrose-notation}
% Keywords: tensor as a node, legs = indices, contraction = joined legs,
% scalars/vectors/matrices/tensors as diagrams, index summation convention,
% trace, identity/copy tensor, diagrammatic vs. index notation.

Index expressions with five or six summations are correct and unreadable.
Diagrams are the working notation of the field for a simple reason: they carry
exactly the same information, and they make the \emph{structure} of a
computation something you can see.
\margintext{Diagrammatic tensor notation is due to Penrose (1971); it entered
numerical practice through the DMRG community in the 1990s and is now standard
in both communities described in \cref{subsec:literature-survey}.}

\begin{itemize}
  \item A tensor is a \textbf{node}; every index is a \textbf{leg}. The rank of
    a tensor is simply its number of legs.
  \item A scalar has no legs, a vector $v_i$ one, a matrix $M_{ij}$ two, a
    general tensor $T_{ijk}$ three, and so on.
  \item Nothing is lost in the translation --- the diagram carries precisely
    what the index expression carries.
\end{itemize}

There are only two moves, and every construction in this course is assembled
from them: \emph{place nodes} and \emph{join legs}.

\begin{itemize}
  \item Joining two legs means \textbf{summing over that shared index}:
    \begin{equation*}
      \sum_j M_{ij} N_{jk}
      \quad\longleftrightarrow\quad
      \text{two nodes, one bond.}
    \end{equation*}
  \item \textbf{Dangling legs} are the free indices of the result. Counting them
    tells you what object you have built.
  \item A \textbf{closed loop} has no free legs left, so the diagram evaluates to
    a number: it is a trace.
\end{itemize}

A small vocabulary of recurring nodes covers most diagrams you will meet:

\begin{itemize}
  \item \textbf{Identity} --- a bare line.
  \item \textbf{Copy (or $\delta$) tensor} --- a branching node, one index shared
    by three legs.
  \item \textbf{Isometry} --- drawn shaded; contracted with its conjugate it
    gives the identity.
\end{itemize}

\begin{notation}
  Throughout these notes physical legs point \textbf{down} and virtual (bond)
  legs run \textbf{horizontally}. Conventions differ between papers; the
  diagrams only mean something once the convention is fixed.
\end{notation}

\begin{remark}
  The real payoff is not compactness. It is that the \emph{order} of contraction
  is visible in the diagram --- and contraction order is what determines cost.
  Two evaluation orders of the same diagram can differ by many orders of
  magnitude in runtime while producing identical numbers.
\end{remark}

\subsection{Literature Survey}
\label{subsec:literature-survey}
% Keywords: DMRG (White 1992), matrix product states, tensor train
% decomposition (Oseledets), quantics/quantized tensor trains (QTT),
% tensor networks for PDEs, tensor networks for CFD, review articles.

Before the technical material, a short orientation --- partly because the
subject has an unusual history. The same ansatz was invented twice, in two
communities that did not talk to each other, and \emph{both} vocabularies remain
in active use. You will meet both in this course and in the papers we draw on.

\begin{center}
  \begin{tabular}{ll}
    \toprule
    \textsc{Quantum many-body} & \textsc{Numerical analysis} \\
    \midrule
    Matrix product state (MPS)   & Tensor train (TT) \\
    Bond dimension $\chi$        & TT-rank \\
    Entanglement across a cut    & Singular-value decay \\
    DMRG sweep                   & Alternating least squares \\
    \bottomrule
  \end{tabular}
\end{center}

\noindent
The two columns denote the same objects and the same algorithms.

\paragraph{Where it came from.}
\margintext{Note the seventeen-year gap between White (1992) and Oseledets
(2011): the numerical-analysis rediscovery was independent, and only later were
the dictionaries matched up.}
\begin{itemize}
  \item \textbf{DMRG} --- variational optimization over matrix product states,
    born in condensed-matter physics~\cite{white1992}.
  \item \textbf{MPS in modern language} --- the standard reference for canonical
    forms and sweeping~\cite{schollwoeck2011}.
  \item \textbf{Tensor-train decomposition} --- the same ansatz rederived for
    numerical linear algebra, with its own error analysis~\cite{oseledets2011}.
  \item \textbf{Quantics} --- binary indexing gives $O(d \log N)$ approximation
    of $N$-$d$ tensors~\cite{khoromskij2011}.
  \item \textbf{TT-cross} --- build a tensor train from samples, without ever
    forming the full array~\cite{oseledets2010ttcross}.
\end{itemize}

\paragraph{Tensor networks meet fluids.}
\begin{itemize}
  \item \textbf{Turbulence has exploitable structure.} Correlations between
    widely separated scales are weak, so the bond dimension stays
    modest~\cite{gourianov2022}.
  \item \textbf{Kinetic plasmas.} Vlasov--Poisson solved in MPS form with large
    compression~\cite{ye2022vlasov}.
  \item \textbf{Reduced-order models.} Wall-bounded flows represented at a few
    percent of the DNS variable count~\cite{kiffner2023rom}.
  \item \textbf{A full stack.} Incompressible flow around arbitrary immersed
    geometry, end to end~\cite{peddinti2024cfd}.
\end{itemize}

\paragraph{Where the field is now.}
\margintext{Our own route through the subject: the 2024 framework fixed the
encoding and the projection loop; the 2026 turbulence study asked how far the
ranks stretch in 3D; the 2025 space--time paper removed the time loop
altogether. Parts~\partref{3} and~\partref{4} follow that order.}
\begin{itemize}
  \item \textbf{3D turbulence} --- compression, simulation and synthesis at
    $\mathrm{Re}_\lambda = 315$ on a $1024^3$ grid~\cite{pisoni2026turbulence}.
  \item \textbf{Compressible flow} --- a 2D Euler solver, finalist in the
    Airbus--BMW Quantum Computing Challenge 2024~\cite{peddinti2025compressible}.
  \item \textbf{Space--time and data-driven} --- space and time in a single MPS,
    together with MPS-DMD~\cite{peddinti2025spacetime}.
\end{itemize}

Parts~\partref{3} and~\partref{4} are built directly on these last
three.

\subsection{Numerical Methods: A Brief Introduction}
\label{subsec:numerical-methods-intro}
% Keywords: curse of dimensionality, grid-based discretization, finite
% differences, why classical solvers struggle at high resolution/dimension,
% compressed representations, low-rank structure.

The motivation for all of this is a scaling wall, so it is worth stating the
wall carefully before claiming to climb it.

\paragraph{The curse of dimensionality.}
\begin{itemize}
  \item Discretize a field on a mesh with $N = 2^n$ points per dimension: in $d$
    dimensions that is $2^{nd}$ unknowns.
  \item A $1024^3$ mesh is already $\sim\!10^9$ unknowns --- \emph{per field,
    per timestep}.
  \item In CFD the cost driver is resolution, and three-dimensional direct
    numerical simulation needs $\sim\!\mathrm{Re}^{9/4}$ degrees of freedom
    (\cref{subsec:navier-stokes}).
  \item Exascale hardware does not close this gap. It moves it by a couple of
    orders of magnitude.
\end{itemize}

\paragraph{Why structure is the only way out.}
\margintext{Every classical accelerator --- adaptive meshes, multigrid, spectral
methods, POD and reduced-order models --- is an exploitation of structure. The
tensor-network ansatz is another one, not a different kind of thing.}
\begin{itemize}
  \item The standard toolkit --- finite differences, volumes or elements,
    explicit or implicit stepping, a linear solve per step --- does not change
    the scaling. It changes the constant.
  \item \textbf{The tensor-network bet:} the exponentially long vector is not an
    arbitrary point in $\mathbb{R}^{2^n}$. A smooth, multiscale field lives on a
    low-rank manifold described by $O(n\chi^2)$ numbers.
  \item \textbf{The price:} exact linear algebra is replaced by approximate
    arithmetic, with a controlled truncation at every single step.
\end{itemize}

That trade --- exponential memory for a truncation error you must monitor --- is
the recurring theme of the entire course.

\subsection{Matrix Product States}
\label{subsec:mps}
% Keywords: MPS/tensor train ansatz, bond dimension, canonical form
% (left/right/mixed), SVD-based compression, singular value decay,
% area law of entanglement, exact vs. approximate ranks, orthogonality center.

\begin{definition}[Matrix product state / tensor train]
  A vector with $2^n$ components is written as a chain of rank-3 tensors,
  \begin{equation}
    v_{i_1 i_2 \cdots i_n}
      = A^{(i_1)}_1 A^{(i_2)}_2 \cdots A^{(i_n)}_n ,
    \label{eq:mps}
  \end{equation}
  where each $A^{(i_k)}_k$ is a matrix of size at most $\chi \times \chi$. The
  \textbf{bond dimension} $\chi$ is the size of the virtual legs.
\end{definition}

The parameter count of~\eqref{eq:mps} is at most $2n\chi^2$ rather than $2^n$.
That is exponential compression --- \emph{provided $\chi$ stays bounded as $n$
grows}. Essentially everything that follows is about when that proviso holds.

\paragraph{Why an MPS can work at all.}
\begin{itemize}
  \item Cut the chain at bond $k$. The bond dimension needed there is the rank of
    $v$ reshaped into a matrix across that cut.
  \item That rank is governed by \textbf{singular-value decay} --- which
    physicists read as the entanglement between the two halves.
  \item For ground states of local Hamiltonians an \textbf{area law} bounds the
    entanglement independently of system size, hence $\chi$ independent of $n$.
  \item Smooth, multiscale functions behave the same way. Rough, noisy or
    randomly structured data do not.
\end{itemize}

\paragraph{Building an MPS: TT-SVD.}
\begin{itemize}
  \item Reshape the vector, take an SVD, split off one tensor, and repeat along
    the chain.
  \item Truncating small singular values at each cut is
    \textbf{quasi-optimal}: within a factor $\sqrt{n}$ of the best rank-$\chi$
    approximation~\cite{oseledets2011}.
  \item Exact and fully controlled --- but it needs the dense $2^n$ vector in
    memory, so it is practical only for modest $n$.
\end{itemize}

\Cref{subsec:finding-qtt} returns to this with methods that build an MPS
\emph{without} ever forming the full array; that is what makes the whole
programme viable.

\paragraph{Canonical forms and rounding.}
\margintext{Gauge freedom is not a nuisance to be tolerated; fixing it is what
makes truncation optimal and inner products cheap. In practice one carries the
orthogonality centre along with the state.}
\begin{itemize}
  \item The factorization~\eqref{eq:mps} is not unique: insert $G G^{-1}$ on any
    bond and nothing changes. Fixing this \textbf{gauge freedom} gives the
    canonical forms.
  \item Sweeping QR from the left and from the right produces \textbf{left,
    right} and \textbf{mixed} canonical forms, the last carrying an
    orthogonality centre.
  \item \textbf{Rounding is the workhorse of the whole subject.} Every operation
    inflates ranks --- a product of two MPSs has bond dimension $\chi_1\chi_2$
    --- so one recompresses back to $\chi$ once per step, every step.
\end{itemize}

\subsection{Quantics Tensor Trains: Functions and Operators}
\label{subsec:qtt-functions-operators}
% Keywords: binary/quantics encoding of a grid, QTT representation of
% functions, exponential compression, QTT rank vs. function smoothness,
% representing differential/difference operators in QTT form, QTT of
% derivative/Laplacian operators.

So far an MPS is a way of factorizing a long vector. The quantics trick is what
turns it into a statement about \emph{physics}, by making each tensor in the
chain correspond to a length scale.

\begin{definition}[Quantics encoding]
  Index a grid point by its binary digits,
  \begin{equation}
    x = \sum_{k=1}^{n} i_k 2^{-k}, \qquad i_k \in \{0,1\}.
    \label{eq:quantics}
  \end{equation}
  A length-$2^n$ vector is thereby reshaped into an $n$-leg tensor, which is then
  written as an MPS. The result is a \textbf{quantics tensor train} (QTT).
\end{definition}

\begin{itemize}
  \item Each tensor now carries \textbf{one length scale}: leg~1 the coarsest,
    leg~$n$ the finest.
  \item The chain becomes a multiscale, wavelet-like object, and the bond between
    legs $k$ and $k+1$ measures the coupling \emph{between scales}.
\end{itemize}

\paragraph{Which functions are low-rank?}
\begin{itemize}
  \item Exact and tiny: $e^{ax}$, $\sin$ and $\cos$ all have $\chi = 2$.
  \item Polynomials of degree $p$ have $\chi = p+1$.
  \item Sums and products of these stay low-rank: ranks add and multiply, and
    then you round.
  \item Doubling the resolution costs $O(n)$ additional parameters rather than
    $O(2^n)$.\margintext{The rule of thumb worth memorizing: rank tracks
    smoothness and scale separation, \emph{not} resolution.}
\end{itemize}

\paragraph{Bit ordering in two and three dimensions.}
With more than one coordinate there is a genuine choice to make:

\begin{itemize}
  \item \textbf{Stacked} --- $x_1 x_2 \cdots x_n\, y_1 y_2 \cdots y_n$: all $x$
    scales, then all $y$ scales.
  \item \textbf{Interleaved} --- $x_1 y_1 x_2 y_2 \cdots$: the coarsest scale of
    both coordinates, then the next, and so on.
\end{itemize}

Interleaving keeps each \emph{scale} local along the chain, which matches the
physics of a turbulent cascade. This is not a cosmetic choice: it decides
whether non-Gaussian turbulence statistics survive
compression~\cite{pisoni2026turbulence}. We return to the evidence in
\cref{subsec:expressivity}.

\paragraph{Operators as MPOs.}
\begin{definition}[Matrix product operator]
  The operator analogue of an MPS: a chain in which each tensor carries two
  physical legs (one in, one out) in addition to its two virtual bonds.
\end{definition}

\begin{itemize}
  \item Finite-difference shift, first-derivative and Laplacian operators have
    \textbf{exact and tiny} MPO ranks, $\chi_{\mathrm{MPO}} \approx 2\text{--}4$.
  \item Applying an MPO to an MPS costs
    $O(n \chi^2 \chi_{\mathrm{MPO}}^2)$, plus a rounding step.
  \item The dense $2^n \times 2^n$ matrix is never formed --- it never exists in
    memory at any point.
\end{itemize}

\subsection{Solving the Heat Equation (Explicit Euler Method)}
\label{subsec:heat-explicit-euler}
% Keywords: 1D/2D heat equation, explicit time-stepping, stability
% (CFL-type) condition, QTT representation of the state, applying the
% discrete Laplacian in QTT form, rounding/recompression after each step,
% numerical example.

We now have enough to solve a PDE. The heat equation is deliberately the easiest
possible choice --- diffusion smooths, so the ranks ought to behave --- which
makes it the right place to see the machinery work and to isolate what can go
wrong.

\paragraph{The setup.}
Take $\partial_t u = \alpha\, \partial_x^2 u$ on $[0,1]$, discretized with
$N = 2^n$ points. Three ingredients suffice, and there is nothing else:

\begin{itemize}
  \item the initial condition $u^0$ as an MPS;
  \item the discrete Laplacian $L$ as an MPO;
  \item one scalar, $\alpha \Delta t / \Delta x^2$.
\end{itemize}

\paragraph{One timestep.}
Explicit Euler is a single line of tensor arithmetic,
\begin{equation}
  u^{k+1} = \left(\mathbb{1} + \alpha \Delta t\, L\right) u^{k},
  \label{eq:heat-explicit}
\end{equation}
consisting of exactly two operations: an \textbf{MPO--MPS product}, and then
\textbf{rounding} to a tolerance $\varepsilon$ (or to a fixed $\chi$).

\begin{itemize}
  \item Cost is $O(n\chi^3)$ per step with $O(n\chi^2)$ memory --- both
    \textbf{logarithmic in the mesh size}.
  \item Skipping the rounding is the classic first mistake: ranks double at every
    step, and within twenty steps you are back to storing $2^n$ numbers.
\end{itemize}

\paragraph{What to watch.}
\margintext{Two error sources, kept apart: discretization
$O(\Delta t, \Delta x^2)$ and truncation. Monitor both --- and always know which
one currently dominates.}
\begin{itemize}
  \item \textbf{Compression does not rescue stability.} The condition is still
    $\alpha \Delta t / \Delta x^2 \le 1/2$, and with $\Delta x = 2^{-n}$ that is
    punishing. This is precisely the motivation for
    \cref{subsec:heat-implicit-euler}.
  \item \textbf{Watch $\chi(t)$.} Diffusion smooths the solution, so the bond
    dimension should stay flat or fall. If it climbs, something is wrong with
    the encoding or with the rounding tolerance.
\end{itemize}

\begin{exercise}
  Run the companion notebook: explicit Euler at $n = 20$, i.e.\ $10^6$ grid
  points, on a laptop. Then remove the rounding step and observe how many
  timesteps you get before memory runs out.
\end{exercise}

\section{Part II: Linear Algebra with Tensor Networks}
\label{sec:part2}
\lecture{29}{09}{2026}

% Mirrors slides.qmd "Part II" outline. One subsection per outline item.

Part~\partref{1} carried out every operation with two primitives: apply an
operator, then round. That is enough for an explicit scheme and nothing else.
The moment a method requires \emph{solving} rather than \emph{applying}, we need
genuine linear algebra on the MPS manifold.

Part~\partref{2} supplies the two missing capabilities and spends them
immediately. First, solving $Ax = b$ in tensor format
(\cref{subsec:linear-solvers}), which buys unconditional stability for the heat
equation (\cref{subsec:heat-implicit-euler}) and returns in Part~\partref{3} as
the pressure-Poisson step. Second, constructing an MPS from a function we can
only sample (\cref{subsec:finding-qtt}), which is what lets us put an aircraft
wing on a $2^{30}$-cell mesh without touching $2^{30}$ cells
(\cref{subsec:tt-cross-geometry}). \Cref{subsec:part2-open-problems} records
what remains unsolved.

\subsection{Linear System Solvers}
\label{subsec:linear-solvers}
% Keywords: solving $Ax=b$ in TT/QTT format, ALS (alternating least
% squares), DMRG-style two-site sweeps, AMEn (alternating minimal energy),
% conjugate gradient / GMRES in tensor format, convergence criteria,
% preconditioning.

The task: $A$ is given as an MPO, $b$ as an MPS, and we want $x$ as an MPS of
bond dimension $\chi$.

\begin{itemize}
  \item Direct inversion is meaningless here. $A^{-1}$ has no reason to be
    low-rank even when $A$ and $x$ both are. (The QTT Laplacian is the
    instructive exception: its inverse \emph{does} admit an explicit low-rank
    representation~\cite{kazeev2012laplace}, which is precisely why that case is
    remarked upon.)
  \item Instead \textbf{recast the problem as a minimization over the MPS
    manifold},
    \begin{equation}
      \min_x \norm{Ax - b}^2
      \qquad\text{or}\qquad
      \min_x \tfrac{1}{2}\, x^{\top}\! A x - x^{\top} b ,
      \label{eq:tt-linear-min}
    \end{equation}%
    the second form being available when $A \succ 0$, i.e.\ symmetric positive
    definite. Prefer it when you can: it avoids squaring the condition number,
    which the residual form does.
  \item The manifold of fixed-rank tensor trains is not a vector space, so we
    cannot optimize all tensors at once. We optimize \textbf{one tensor at a
    time}, which is a linear problem in that tensor alone.
\end{itemize}

\paragraph{Sweeping solvers: ALS and DMRG.}
\margintext{The one-site/two-site distinction is the same one that separates
White's original DMRG from its later single-site variants. The trade is rank
adaptivity against local cost.}
\begin{itemize}
  \item \textbf{ALS, or one-site DMRG.} Freeze every tensor but one. With the
    chain in mixed canonical form centred on that tensor, the problem collapses
    to a small dense linear system of size $\chi^2 d$, solved exactly. Sweep
    left to right and back~\cite{holtz2012als}.
  \item Each local solve is exact, so the global objective decreases
    monotonically. Convergence to a global optimum is not guaranteed; ALS can
    stall in a local minimum of the fixed-rank manifold.
  \item Its limitation: $\chi$ is fixed in advance and never adapts.
  \item \textbf{Two-site DMRG.} Optimize a \emph{pair} of neighbouring tensors,
    then split them by SVD. The truncation tolerance in that SVD sets the new
    bond dimension, so $\chi$ adapts on its own. The local problem is now of
    size $\chi^2 d^2$ and the local cost rises to $O(\chi^6)$.
  \item Two-site updates also escape local minima that trap one-site
    sweeps, because the pair carries variational freedom that neither tensor has
    alone.
\end{itemize}

\paragraph{AMEn and Krylov alternatives.}
\begin{itemize}
  \item \textbf{AMEn}~\cite{dolgov2014amen} --- one-site sweeps enriched with
    directions taken from the current residual $b - Ax$. This recovers rank
    adaptivity at roughly one-site cost, and is often the best default.
  \item \textbf{Krylov methods in TT format} --- CG or GMRES with a rounding
    step after every operation. Easy to build on top of existing TT arithmetic.
  \item But rank growth per iteration is severe: each matrix--vector product
    multiplies $\chi$ by $\chi_{\mathrm{MPO}}$ and each vector addition adds
    ranks, so truncation must be aggressive, and aggressive truncation destroys
    the Krylov subspace's orthogonality.
  \item Preconditioning in TT form remains open
    (\cref{subsec:part2-open-problems}).
\end{itemize}

\begin{remark}[Rule of thumb]
  Use sweeping methods for structured PDE operators, where the MPO is known in
  closed form and warm starts are available. Use Krylov methods when the
  operator is effectively a black box.
\end{remark}

\paragraph{What the solve actually costs.}
This is not a peripheral concern. In our incompressible CFD solver the
DMRG-type Poisson solve dominates every other part of the
pipeline~\cite{peddinti2024cfd}:

\begin{center}
  \begin{tabular}{ll}
    \toprule
    \textsc{Quantity} & \textsc{Value} \\
    \midrule
    Share of total runtime  & $80.1\%$ \\
    Worst-case local cost   & $O(n\chi^6)$ \\
    Observed total cost     & $O(n\chi^{4.1})$ \\
    \bottomrule
  \end{tabular}
\end{center}

\noindent
\margintext{The measured exponent $4.1$ against a worst case of $6$ is where all
the practical engineering lives. Warm-starting from the previous timestep is the
single largest contributor.}
The gap between the worst case and what is measured is where the practical knobs
live: sweep count, local tolerance, rank cap, and above all
\textbf{warm-starting from the previous timestep}. The worst-case $\chi^6$
assumes the local eigen- or linear solve runs to full accuracy at maximum rank
on every site of every sweep; warm starts mean it almost never does.

\subsection{Solving the Heat Equation (Implicit Euler Method)}
\label{subsec:heat-implicit-euler}
% Keywords: implicit time-stepping, unconditional stability, linear
% system per time step, QTT operator inversion, comparison with explicit
% scheme, cost per step vs. accuracy/stability trade-off.

Return to the heat equation~\eqref{eq:heat} of
\cref{subsec:heat-explicit-euler}, now with backward Euler:
\begin{equation}
  \left(\mathbb{1} - \alpha \Delta t\, L\right) u^{k+1} = u^{k}.
  \label{eq:heat-implicit}
\end{equation}

\paragraph{Why bother.}
\begin{itemize}
  \item Equation~\eqref{eq:heat-implicit} requires a \textbf{linear solve} every
    step instead of a matrix--vector product: strictly more work per step.
  \item The payoff is unconditional stability. The amplification factor is
    $(1 + \alpha \Delta t\, \lambda)^{-1}$ for every eigenvalue $\lambda \ge 0$
    of $-L$, so it never exceeds one whatever $\Delta t$ is. The step size is
    then set by \textbf{accuracy}, not by $\Delta x^2$.
  \item With $\Delta x = 2^{-n}$ and $n \gtrsim 20$, the explicit bound
    $\Delta t \le \Delta x^2/2\alpha$ is unusable. That is the entire argument,
    and it strengthens with every bit of resolution quantics buys.
\end{itemize}

\begin{remark}
  Note the shape of the trade. Quantics makes very large $n$ affordable in
  memory; the explicit stability condition then makes large $n$ unaffordable in
  time. The compression only pays off in combination with an implicit scheme,
  which is only affordable because of \cref{subsec:linear-solvers}.
\end{remark}

\paragraph{Assembling and solving.}
\begin{itemize}
  \item Build $A = \mathbb{1} - \alpha \Delta t\, L$ directly as an MPO. Its
    rank is $\approx 3$ and the cores are written down
    analytically~\cite{kazeev2012laplace}; there is no numerical construction
    step and no assembly error.
  \item The right-hand side is simply the previous state MPS.
  \item Solve with DMRG or AMEn, \textbf{warm-started from $u^k$}. Consecutive
    solutions differ by $O(\Delta t)$, so one or two sweeps usually suffice.
\end{itemize}

\begin{remark}[How to compare schemes]
  Compare \textbf{total wall-clock time to a target error}, never cost per step.
  Implicit stepping loses on the second measure and wins on the first, because
  it takes far fewer steps.
\end{remark}

\paragraph{The error budget.}
The budget now has three parts rather than two:

\begin{enumerate}
  \item time discretization, $O(\Delta t)$;
  \item space discretization, $O(\Delta x^2)$;
  \item \textbf{the solver residual} $\norm{Au^{k+1} - u^{k}}$ --- new, and
    entirely under your control.
\end{enumerate}

\noindent
Keep the residual well below the two discretization errors. Otherwise you
converge efficiently to the wrong answer, which is the least useful of all
possible outcomes. A practical rule: target a relative residual one order of
magnitude below the expected discretization error, and no tighter, since
sweeping past that point is pure cost.
\margintext{Forward pointer: this exact solve reappears in
\cref{subsec:qtt-cfd-frameworks} as the pressure-Poisson step of Chorin
projection, where it consumes $80\%$ of the runtime.}

\subsection{How to Find the QTT of Your Function?}
\label{subsec:finding-qtt}
% Keywords: constructing QTT from analytic/black-box functions,
% interpolation vs. exact construction, rank growth with resolution,
% sampling on a quantics grid, adaptive rank truncation.

We have been assuming the initial condition arrives already in MPS form. That
assumption has to be discharged, and how it is discharged decides whether the
method scales. There are three routes.

\begin{center}
  \begin{tabular}{@{}lll@{}}
    \toprule
    \textsc{Route} & \textsc{Cost} & \textsc{Requires} \\
    \midrule
    Analytic cores      & free                        & closed form for $f$ \\
    TT-SVD              & $O(2^n)$                    & the dense vector \\
    TT-cross / TCI      & $O(n\chi^2)$ evaluations    & only a black box $f$ \\
    \bottomrule
  \end{tabular}
\end{center}

\begin{itemize}
  \item \textbf{Analytic.} For the families tabulated in
    \cref{subsec:qtt-functions-operators} the small cores are written down by
    hand~\cite{oseledets2013functions}. Exact: no sampling, no error.
  \item \textbf{From a full array (TT-SVD).} Exact and quasi-optimal
    (\cref{subsec:mps}), but it needs the dense $2^n$ vector, so it is viable
    only for small $n$.
  \item \textbf{Black-box (TT-cross, or TCI).} Build the MPS from $O(n\chi^2)$
    \textbf{evaluations} of $f$, never forming the full
    tensor~\cite{oseledets2010ttcross}.
\end{itemize}

The third route is the one that scales, and the one that makes the method
practical rather than merely elegant.
\margintext{TT-cross is what keeps the programme from being circular. Without
it you would need the $2^n$-vector in order to compress the $2^n$-vector.}

\paragraph{Cross interpolation and maxvol.}
\begin{itemize}
  \item Cross interpolation is a \textbf{skeleton (CUR) decomposition applied
    recursively} down the chain.
  \item A rank-$r$ matrix is exactly reconstructible from $r$ well-chosen rows
    and columns. The only question is \emph{which} ones.
  \item \textbf{maxvol}: choose the $r \times r$ submatrix of maximum volume,
    i.e.\ largest $\abs{\det}$. This is near-optimal and computable greedily by
    row swaps.
  \item Quantics combined with cross interpolation gives \textbf{QTCI}:
    parsimonious representations of multivariate functions at very high
    resolution~\cite{ritter2024qtci}. Mature implementations exist in
    \texttt{xfac} and \texttt{Tensor\-Cross\-Inter\-polation.jl}~\cite{nunez2025tci}.
\end{itemize}

The algorithm is spelled out in \cref{subsec:appendix-tt-cross}.

\paragraph{TT-cross in practice.}
The working recipe is short:

\begin{enumerate}
  \item choose the bit ordering, interleaved or stacked
    (\cref{subsec:qtt-functions-operators});
  \item set a tolerance;
  \item sweep until the cross error plateaus;
  \item round.
\end{enumerate}

\begin{remark}[Failure modes to expect]
  \begin{itemize}
    \item A discontinuous $f$ causes rank blow-up: a step edge has singular
      values decaying like $1/k$, so no finite $\chi$ captures it well.
    \item Poorly placed pivots cause features to be missed entirely. The method
      never samples them, so it never learns they exist, and the reported cross
      error stays small while the answer is wrong.
    \item A noisy $f$ is interpolated faithfully, noise included. Cross
      interpolation is exact on the fibres it samples, which is a virtue for
      clean data and a liability for dirty data.
  \end{itemize}
\end{remark}

\subsection{Encoding Geometry via TT-Cross}
\label{subsec:tt-cross-geometry}
% Keywords: TT-cross / tensor cross interpolation (TCI), maxvol algorithm,
% black-box tensor construction without full evaluation, representing
% complex domains/boundaries in QTT, indicator functions, adaptive
% sampling cost.

Now the payoff, and the first place where the tensor-network constraints force a
genuine modelling decision rather than an implementation detail.

\paragraph{The problem with sharp boundaries.}
\begin{itemize}
  \item The goal is to place an arbitrary immersed object --- a cylinder, an
    airfoil --- on a $2^n$ mesh \textbf{without touching $2^n$ cells}.
  \item The natural object is an indicator function $\theta(q(x))$: one inside
    the body, zero outside, where $q$ is a level-set function vanishing on the
    boundary.
  \item That is exactly the wrong thing for a tensor network. A step edge has
    \textbf{slowly decaying singular values}, so $\chi$ explodes. This is the
    failure mode of \cref{subsec:mps} in its purest form: the indicator is
    maximally non-smooth exactly where the physics happens.
\end{itemize}

The fix is to give up the sharp edge \emph{deliberately}, and then bound what
that costs.

\paragraph{Smoothed masks.}
\margintext{This construction is from our 2024 framework. The point of the
$\sqrt{\ell_0^{D-2}/\alpha}$ bound is that the approximation is a modelling
choice with an error bar, not an undocumented hack.}
Replace the indicator by a smooth mask, with $q$ vanishing on the object
boundary $\partial\Omega$~\cite{peddinti2024cfd}:
\begin{equation}
  m(x) = 1 - e^{-\alpha\left(q(x) + \abs{q(x)}\right)},
  \qquad q|_{\partial\Omega} = 0 .
  \label{eq:smoothed-mask}
\end{equation}

\begin{itemize}
  \item The combination $q + \abs{q}$ vanishes identically wherever $q < 0$, so
    $m$ is exactly zero inside the body and rises smoothly outside it. No
    approximation is made in the interior.
  \item $\alpha$ trades sharpness against rank: large $\alpha$ approaches the
    true indicator and costs bond dimension.
  \item The $\ell_2$ distance to the true indicator scales as
    $\sqrt{\ell_0^{D-2}/\alpha}$ in $D$ dimensions, where $\ell_0$ is the
    boundary length scale. The error is therefore controlled by a single knob
    and decays as $\alpha^{-1/2}$.
\end{itemize}

\paragraph{Putting it together.}
\begin{itemize}
  \item Build $m$ as an MPS by TT-cross to tolerance $\varepsilon$, at a cost
    independent of the mesh size.
  \item Apply it by an \textbf{elementwise (Hadamard) MPS product} on the
    velocity field each step. No-slip is enforced without ever meshing the body,
    and without an immersed-boundary force term.
  \item Inlet and outlet conditions fold into the differential MPOs via
    \textbf{ghost cells}, so open boundary conditions keep the domain small.
  \item Demonstrated on a \textbf{NACA 0040 airfoil}: $\chi = 45$, $18{,}800$
    parameters, a compression factor of $27.9$.
\end{itemize}

\subsection{Open Problems}
\label{subsec:part2-open-problems}
% Keywords: rank blow-up for non-smooth data, robustness of TT-cross,
% error control, nonlinear operators in TT format, scalability to 3D/time.

Worth being explicit about what is not settled. The honest list is short enough
to state and long enough to matter.

\paragraph{Ranks and nonlinearity.}
\begin{itemize}
  \item \textbf{Rank blow-up.} Shocks, sharp interfaces, non-smooth data and
    high-$\mathrm{Re}$ turbulence all push $\chi$ up. Where exactly does the
    method stop paying for itself?
  \item \textbf{Nonlinearity.} Products of MPSs square the bond dimension.
    Picard iteration converges linearly and Newton quadratically, but every
    iteration re-rounds, and the interaction between iteration error and
    truncation error is not understood.
  \item \textbf{Error control.} There is no sharp a-posteriori bound linking the
    per-step truncation tolerance $\varepsilon$ to the PDE error accumulated
    over many steps. In practice one assumes linear accumulation, $O(K\varepsilon)$
    after $K$ steps, and checks it empirically.
  \item \textbf{Robustness of TT-cross.} Pivot selection is heuristic, with no
    guarantee for adversarial or noisy black boxes.
\end{itemize}

\paragraph{Solvers and structure.}
\margintext{Of everything on this page, preconditioning has the largest
immediate payoff. Recall that the Poisson solve is $80\%$ of runtime.}
\begin{itemize}
  \item \textbf{Preconditioning in TT format} is largely unsolved, and it is the
    single biggest lever on solver cost. The difficulty is structural: a good
    preconditioner approximates $A^{-1}$, and low-rank $A$ rarely has low-rank
    inverse.
  \item \textbf{Conservation.} SVD truncation preserves neither mass nor energy
    nor $\nabla\cdot\mathbf{v} = 0$; it is optimal in the Frobenius norm, which
    knows nothing about the physics. Structure-preserving rounding is wide open.
  \item \textbf{Scaling out.} Three dimensions plus time, index ordering at
    scale, memory-bound implementations, and distributed or GPU sweeps.
  \item \textbf{Fair benchmarking.} Honest comparison against tuned classical
    CFD --- multigrid, spectral methods --- on equal hardware.
\end{itemize}

\section{Part III: The Forward Problem - Solving PDEs}
\label{sec:part3}
\lecture{30}{09}{2026}

% Mirrors slides.qmd "Part III" outline. One subsection per outline item.

Parts~\partref{1} and~\partref{2} assembled a complete toolkit on a problem
chosen for its docility. Part~\partref{3} points that toolkit at the
Navier--Stokes equations, where every weakness of the toolkit is provoked at
once: nonlinearity, a global constraint, and complicated geometry.

The order of the argument: what the equations demand
(\cref{subsec:navier-stokes}), how a QTT solver is built
(\cref{subsec:qtt-cfd-frameworks}), why fluids are the right target for this
ansatz rather than an arbitrary one (\cref{subsec:why-fluid-dynamics}), and then
the question everything hinges on --- when does $\chi$ actually stay small
(\cref{subsec:expressivity})?

\subsection{Navier-Stokes Equations}
\label{subsec:navier-stokes}
% Keywords: incompressible/compressible Navier--Stokes, momentum and
% continuity equations, boundary conditions, nondimensionalization,
% Reynolds number, why it is a hard high-dimensional nonlinear PDE.

The incompressible Navier--Stokes equations read
\begin{equation}
  \partial_t \mathbf{v} + (\mathbf{v}\cdot\nabla)\mathbf{v}
    = -\frac{1}{\rho}\nabla p + \nu \nabla^2 \mathbf{v},
  \qquad
  \nabla\cdot\mathbf{v} = 0 .
  \label{eq:navier-stokes}
\end{equation}

They present three difficulties, and --- this is the useful observation --- each
is awkward for a tensor network in a \emph{different} way:

\begin{enumerate}
  \item \textbf{Nonlinear advection}, $(\mathbf{v}\cdot\nabla)\mathbf{v}$:
    products of MPSs, so ranks multiply.
  \item \textbf{A nonlocal pressure constraint}: a global elliptic solve at
    every step.
  \item \textbf{Boundary conditions on complex geometry}: curved interfaces
    couple every length scale at once, hence high rank
    (\cref{subsec:tt-cross-geometry} --- and note that it is the curvature, not
    the sharpness, that costs).
\end{enumerate}

\noindent
\Cref{subsec:tt-cross-geometry} disposed of the third. The first is managed by
rounding after every product. The second is the expensive one.

\paragraph{Why pressure is the hard part.}
\margintext{This is not a tensor-network problem. The globally coupled elliptic
solve dominates runtime in classical incompressible CFD too, which is what makes
the comparison in \cref{subsec:qtt-cfd-results} a fair one.}
\begin{itemize}
  \item Pressure is \textbf{not a dynamical variable}. There is no evolution
    equation for it in~\eqref{eq:navier-stokes}.
  \item It is whatever instantaneously enforces $\nabla\cdot\mathbf{v} = 0$: a
    constraint, not a dynamics.
  \item Taking the divergence of the momentum equation and using
    $\nabla\cdot\mathbf{v} = 0$ to kill the time-derivative and viscous terms
    leaves an elliptic \textbf{Poisson problem},
    \begin{equation}
      \nabla^2 p = -\rho\, \nabla\cdot\!\left[(\mathbf{v}\cdot\nabla)\mathbf{v}\right],
      \label{eq:pressure-poisson}
    \end{equation}%
    whose right-hand side is the only place the nonlinearity enters.
  \item Equation~\eqref{eq:pressure-poisson} couples the whole domain at once:
    the Green's function of $\nabla^2$ has infinite support, so information
    propagates across the mesh within a single step. Every timestep therefore
    contains a globally coupled linear solve. That is where the runtime goes ---
    in classical CFD and in tensor networks alike.
\end{itemize}

\paragraph{Reynolds number and the cascade.}
\begin{itemize}
  \item Nondimensionalize~\eqref{eq:navier-stokes} and exactly one control
    parameter survives: $\mathrm{Re} = UL/\nu$, the ratio of inertial to viscous
    forces.
  \item \textbf{Kolmogorov picture:} energy is injected at the large scale $L$
    and dissipated at the Kolmogorov scale
    $\eta \sim L\,\mathrm{Re}^{-3/4}$, cascading through the scales between.
  \item Resolving both ends in three dimensions therefore needs
    $(L/\eta)^3 \sim \mathrm{Re}^{9/4}$ degrees of freedom --- the estimate
    quoted without justification in \cref{subsec:numerical-methods-intro}.
  \item A commercial aircraft at cruise sits around $\mathrm{Re} \sim 10^7$,
    giving $\sim\!10^{16}$ degrees of freedom. That is the number that makes
    direct numerical simulation impossible.
\end{itemize}

\paragraph{Benchmarks used in these notes.}
\begin{center}
  \begin{tabular}{@{}lp{0.52\textwidth}@{}}
    \toprule
    \textsc{Case} & \textsc{What it tests} \\
    \midrule
    Lid-driven cavity     & closed domain, no immersed geometry; the standard
                            sanity check \\
    Flow past a cylinder  & steady vortex pair, then a K\'arm\'an street;
                            well tabulated \\
    NACA 0040 airfoil     & genuinely non-trivial immersed geometry \\
    Decaying isotropic
    turbulence            & rank behaviour at high $\mathrm{Re}$; the hard
                            case \\
    \bottomrule
  \end{tabular}
\end{center}

\noindent
The compressible cousin of~\eqref{eq:navier-stokes} --- the Euler equations,
with shocks and a nonlinear flux --- is targeted
separately~\cite{Peddinti2025}.

\subsection{QTT Frameworks for CFD}
\label{subsec:qtt-cfd-frameworks}
% Keywords: QTT discretization of velocity/pressure fields, nonlinear
% advection term in TT format, projection/pressure-correction methods,
% lattice Boltzmann in QTT, existing software/frameworks, benchmark flows
% (lid-driven cavity, Burgers' equation).

\paragraph{The encoding.}
\margintext{Every design decision here is inherited from Part~\partref{1}:
interleaving from \cref{subsec:qtt-functions-operators}, low-rank MPOs from the
same place, rounding from \cref{subsec:mps}.}
\begin{itemize}
  \item Each velocity and pressure component becomes \textbf{one MPS} over
    $n = n_x + n_y \,(+\, n_z)$ quantics bits.
  \item Bits are \textbf{interleaved} across coordinates, so each length scale
    stays local along the chain.
  \item Gradient, divergence and Laplacian are MPOs of rank $2$--$4$, built once
    in closed form and reused throughout~\cite{Kazeev2012}.
  \item Everything downstream is MPO application, Hadamard product, and
    rounding. The solver contains no other operation.
\end{itemize}

\paragraph{Chorin projection.}
The scheme is Chorin's projection method~\cite{Chorin1968} carried out in tensor
form~\cite{Peddinti2024}. Three moves per timestep:

\begin{enumerate}
  \item \textbf{Predictor} --- ignore pressure and advance momentum,
    \begin{equation*}
      \mathbf{v}^{*} = \mathbf{v}^{t}
        + \Delta t \left(-(\mathbf{v}^t\cdot\nabla)\mathbf{v}^t
        + \nu\nabla^2\mathbf{v}^t\right).
    \end{equation*}
    The result is generally not divergence-free.
  \item \textbf{Pressure Poisson} --- the DMRG linear solve of
    \cref{subsec:linear-solvers}, and the expensive step,
    \begin{equation*}
      \nabla^2 p = \frac{\rho}{\Delta t}\, \nabla\cdot\mathbf{v}^{*}.
    \end{equation*}
  \item \textbf{Projection} --- subtract the pressure gradient to land back on
    the divergence-free manifold,
    \begin{equation*}
      \mathbf{v}^{t+1} = \mathbf{v}^{*} - \frac{\Delta t}{\rho}\nabla p .
    \end{equation*}
\end{enumerate}

\begin{remark}
  Steps 2 and 3 are a Helmholtz--Hodge decomposition: any vector field splits
  uniquely into a divergence-free part and a gradient, and the Poisson solve is
  how you find the gradient part in order to discard it.
\end{remark}

\begin{remark}[Which Laplacian goes in step 2]
  \margintext{Costs a day to find, ten seconds to fix, and looks exactly like
  \enquote{the solver is a bit inaccurate} the whole time.}
  The operator in step~2 must be $\nabla\cdot\nabla$ \emph{assembled from the
  same difference operators} used for the gradient in step~3 and the divergence
  in step~2. Substitute the compact three-point Laplacian --- which is the
  obvious thing to reach for, since \cref{subsec:qtt-functions-operators} built
  it and it has a lower MPO rank --- and the projection no longer projects.

  The algebra is one line. With $D$ the discrete gradient,
  $\nabla\cdot\mathbf{v}^{t+1} = D\cdot\mathbf{v}^{*} - \Delta t\,(D\cdot D)p$,
  which vanishes identically only if the operator inverted in step~2 was
  $D\cdot D$. With centred differences $D\cdot D \ne L$, and the residual
  divergence is $O(1)$ rather than $O(\text{solver tolerance})$. Measured on one
  step of our demonstrator: $\norm{\nabla\cdot\mathbf{v}}$ falls from $7.7$ to
  $4\times10^{-3}$ when the operators match, and by a factor of a few when they
  do not.
\end{remark}

\paragraph{Geometry and readout.}
\begin{itemize}
  \item \textbf{Geometry.} The smoothed MPS mask of
    \cref{subsec:tt-cross-geometry} is applied elementwise after the predictor;
    ghost-cell inlet and outlet conditions live inside the MPOs.
  \item \textbf{Self-consistent readout} --- the part that makes the method
    usable rather than merely fast:
    \begin{itemize}
      \item a coarse-grained field costs $O(n\chi^3)$, by contracting the fine
        legs against uniform vectors;
      \item a single grid value costs $O(\chi^3)$, by contracting the chain
        against the binary digits of that index.
    \end{itemize}
  \item The dense $2^n$ vector is \textbf{never reconstructed}, at any point in
    the pipeline.
\end{itemize}

\begin{remark}
  Without a compressed readout the whole advantage evaporates: you would
  compress the computation and then decompress it in order to look at it. This
  is a general lesson about compressed solvers, not a detail of this one.
\end{remark}

\paragraph{Results and scaling.}
\label{subsec:qtt-cfd-results}
Incompressible flow past a cylinder~\cite{Peddinti2024}:

\begin{center}
  \begin{tabular}{ll}
    \toprule
    \textsc{Quantity} & \textsc{Value} \\
    \midrule
    Mesh                        & $2^{30}$ cells \\
    Bond dimension              & $\chi = 30$ \\
    Parameters                  & $24{,}200$ \\
    Compression                 & $\sim\!4.4\times 10^{4}$ \\
    Reynolds range (all cases)  & $1.7$--$2230$ \\
    Observed runtime            & $\sim n\chi^{4.1}$ \\
    \bottomrule
  \end{tabular}
\end{center}

\noindent
Runtime is linear in $n$ and therefore \textbf{poly-logarithmic in the mesh
size} $2^n$. The results were validated against Ansys Fluent on the squared
geometries.

\begin{remark}[What the teaching demonstrator reproduces, and what it does not]
  The companion solver in \texttt{notebooks/qtade\_cfd.py} is a few hundred
  lines of NumPy with homogeneous Dirichlet conditions on every wall, no
  ghost-cell inlet or outlet, and no immersed geometry. It is not the code that
  produced the table above, and it is not validated against anything.

  What it does reproduce is the \emph{cost structure}, which is the part worth
  carrying away. On a $64\times64$ grid over 40 timesteps the pressure solve
  takes $82\%$ of the runtime, against the $80.1\%$ reported for the full
  solver. Scaling from $n=5$ to $n=8$ multiplies the cell count by 64 and the
  runtime by about 20 --- sublinear in cells, and well above the ideal $O(n)$,
  the excess being interpreter overhead and local conjugate-gradient iterations
  rather than anything about tensor networks. Take the shape of that curve from
  the demonstrator and the constant from the paper.
\end{remark}

\begin{remark}[Truncation does not respect the constraint]
  SVD truncation is optimal in the Frobenius norm, and the Frobenius norm knows
  nothing about $\nabla\cdot\mathbf{v}=0$. Every rounding step puts a little
  divergence back and the next projection removes it again; the balance is set
  by the rank cap. Measured after 20 steps of the demonstrator:

  \begin{center}
    \begin{tabular}{@{}lr@{}}
      \toprule
      $\chi_{\max}$ & $\norm{\nabla\cdot\mathbf{v}}$ \\
      \midrule
      8  & $1.2\times10^{1}$ \\
      16 & $1.6\times10^{0}$ \\
      32 & $4.1\times10^{-2}$ \\
      48 & $5.9\times10^{-4}$ \\
      \bottomrule
    \end{tabular}
  \end{center}

  \noindent
  Structure-preserving rounding --- truncation that conserves mass, energy or
  the constraint --- is open (\cref{subsec:part2-open-problems}). Until it
  exists, $\chi_{\max}$ is the knob and the divergence is the diagnostic you
  watch.
\end{remark}

\paragraph{Other frameworks in the field.}
\begin{itemize}
  \item \textbf{Reduced-order models} for wall-bounded flows, at a few percent
    of the DNS variable count~\cite{Kiffner2023}.
  \item \textbf{Turbulence PDFs} computed directly in tensor-network form,
    rather than by sampling a reconstructed field~\cite{Gourianov2025}.
  \item \textbf{Vlasov--Poisson} for kinetic plasmas~\cite{Ye2022}.
  \item \textbf{Lattice Boltzmann} in tensor-train form: a different
    discretization, the same compression idea.
\end{itemize}

The common thread is that the ansatz is portable across discretizations, and
that the global solve is always the hard part.

\subsection{Why Fluid Dynamics?}
\label{subsec:why-fluid-dynamics}
% Keywords: computational cost of high-resolution CFD, turbulence and
% scale separation, industrial/scientific relevance, motivation for
% compressed tensor representations.

Two separate questions, worth keeping apart: why fluids are worth the effort,
and why they happen to suit this particular ansatz.

\paragraph{Fluid dynamics is where the cycles go.}
\begin{itemize}
  \item CFD is a top consumer of HPC worldwide: aerodynamics, weather and
    climate, combustion, biomedical flow.
  \item The bottleneck is resolution, and resolution costs
    $\sim\!\mathrm{Re}^{9/4}$ in three dimensions.
  \item Exascale still cannot simulate a full aircraft at flight Reynolds
    number.
  \item Everything cheaper than DNS --- RANS, LES --- is a \textbf{model}, with
    tuned closures and limited transfer between regimes.
\end{itemize}

\paragraph{Why it suits tensor networks.}
\margintext{Gourianov \emph{et al.} (2022) made this argument first and
quantitatively. It is the observation that turned a plausible analogy into a
research programme.}
\begin{itemize}
  \item Turbulence has precisely the structure the ansatz wants: a
    \textbf{cascade} in which well-separated scales couple
    weakly~\cite{Gourianov2022}.
  \item Quantics provides a scale-by-scale representation for free
    (\cref{subsec:qtt-functions-operators}), so the natural hierarchy of the
    physics \emph{is} the natural hierarchy of the ansatz. Weak inter-scale
    coupling is exactly what a small bond dimension encodes.
  \item It is also an honest stress test: nonlinear, multiscale, geometrically
    complex, with unambiguous quantitative benchmarks.
  \item And it is a bridge to hardware: the same encoding is what a
    fault-tolerant quantum CFD algorithm would consume and return.
\end{itemize}

\subsection{Expressivity}
\label{subsec:expressivity}
% Keywords: what functions/fields admit low QTT rank, smoothness vs.
% rank, turbulent vs. laminar flow compressibility, empirical rank
% studies, limits of expressivity.

Everything so far assumed that $\chi$ stays small. This subsection asks when
that is actually true, and it is the most important one in Part~\partref{3}.

\paragraph{The central question.}
Two sub-questions, with different answers:

\begin{itemize}
  \item How does $\chi$ grow with \textbf{resolution}, at fixed physics?
  \item How does $\chi$ grow with \textbf{Reynolds number}?
\end{itemize}

\noindent
The first is benign --- that is the quantics promise of
\cref{subsec:qtt-functions-operators}. The second is the real question, it has
no theory behind it (\cref{subsec:part3-outlook}), and the answer is empirical.

\paragraph{3D turbulence at $\mathrm{Re}_\lambda = 315$.}
\margintext{Pisoni \emph{et al.} (2026): DNS snapshots on $1024^3 = 2^{30}$
points, encoded as 30-tensor trains. The largest such study to date, and the
source of the numbers below.}
\begin{center}
  \begin{tabular}{@{}llp{0.46\textwidth}@{}}
    \toprule
    $\chi$ & \textsc{Retained} & \textsc{What survives} \\
    \midrule
    $1000$ & $2.2\%$  & the full flat inertial range of $E(k)$ \\
    $100$  & $0.03\%$ & enough to run the solver stably for 9--10 turnover
                        times \\
    \bottomrule
  \end{tabular}
\end{center}

\noindent
That is four orders of magnitude of compression with the energy spectrum
essentially intact~\cite{Pisoni2026}.

\paragraph{What breaks first.}
\begin{itemize}
  \item \textbf{The spectrum is the easy diagnostic.} It still looks fine long
    after the flow does not. $E(k)$ is a second-order two-point statistic, and
    second-order statistics are exactly what an $\ell_2$-optimal truncation
    protects. Reporting only $E(k)$ therefore overstates the quality of a
    compression.
  \item The demanding diagnostics are the \textbf{PDFs of velocity increments}
    and the \textbf{flatness}, i.e.\ the intermittency statistics. These are
    high-order and live in the tails, where truncation error concentrates.
  \item \textbf{Bit ordering decides this.} Stacked ordering over-smooths the
    field at low $\chi$; interleaved ordering preserves the non-Gaussian
    statistics.
  \item But below $\chi \approx 500$ the interleaved encoding injects spurious
    small-scale fluctuations. Both encodings fail --- differently.
\end{itemize}

\begin{remark}[The ordering in which diagnostics recover]
  The claim above is easy to check in two dimensions, and worth checking because
  the ordering is so lopsided. Compressing a synthetic intermittent field
  (a $k^{-5/3}$ spectrum modulated by a lognormal multiplicative cascade, true
  flatness $37.6$):

  \begin{center}
    \begin{tabular}{@{}lrrr@{}}
      \toprule
      $\chi$ & $\ell_2$ error & $E(k)$ error, $k<40$ & flatness \\
      \midrule
      16 & $0.61$ & $37\%$   & $119$ \\
      32 & $0.42$ & $13\%$   & $64$ \\
      64 & $0.21$ & $1.4\%$  & $44$ \\
      \bottomrule
    \end{tabular}
  \end{center}

  \noindent
  At $\chi = 64$ the spectrum is right to $1.4\%$ while the pointwise error is
  still $21\%$ and the flatness is $16\%$ high. The diagnostic that survives
  compression longest is therefore the one least able to tell you that the
  compression was safe. State this to anyone who shows you a compressed
  turbulent field and an $E(k)$ plot.
\end{remark}

\begin{remark}[The honest claim]
  What the evidence supports is that turbulence compresses \emph{well at
  moderate} $\mathrm{Re}$. What it does not support is any claim that turbulence
  is low-rank at arbitrary $\mathrm{Re}$ and at every scale. The $\chi = 100$
  ceiling in the current 3D solver is a memory-bound implementation limit, not a
  physical statement.

  A second, cheaper measurement makes the shape of the dependence visible
  without a DNS: synthesize fields with an inertial range of increasing width
  and compress them to a fixed tolerance. With $k_{\max} = 4, 8, 16, 32, 64$ the
  bond dimension needed at $10^{-2}$ goes $14, 26, 41, 68, 128$. The rank tracks
  the \emph{width of the inertial range}, which is the thing $\mathrm{Re}$
  controls --- not the grid, which is the thing quantics already handles.
\end{remark}

The flip side is genuinely encouraging. Random low-rank tensor trains with
$\chi = 10$, cascaded across scales, \textbf{synthesize} divergence-free
turbulent-like snapshots --- correct spectrum, correct flatness --- at
logarithmic cost~\cite{Pisoni2026}. Whatever the ansatz is missing, it
is not the gross statistical structure of turbulence.

\subsection{Outlook}
\label{subsec:part3-outlook}
% Keywords: scaling to 3D/turbulent regimes, adaptive time stepping,
% coupling with machine learning, open numerical-analysis questions.

\paragraph{Solvers and scale.}
\begin{itemize}
  \item \textbf{Higher $\mathrm{Re}$, full 3D.} Lifting the $\chi = 100$ ceiling
    is an engineering problem before it is a physics one.
  \item \textbf{Better solvers.} TT preconditioning, distributed and GPU sweeps,
    mixed precision. Recall that the Poisson solve is $80\%$ of runtime, so
    nothing else is worth optimizing first.
  \item \textbf{Adaptivity.} Choose $\chi(t)$ and $\Delta t$ from a live error
    estimate rather than a fixed cap; tangent-space integrators supply the
    machinery~\cite{Lubich2015}.
  \item \textbf{Structure-preserving truncation.} Maintain
    $\nabla\cdot\mathbf{v} = 0$ and the conserved quantities under SVD rounding.
\end{itemize}

\paragraph{Physics and theory.}
\margintext{The theory gap is the takeaway for an applied-mathematics audience:
there is at present no sharp rank bound for Navier--Stokes solutions as a
function of $\mathrm{Re}$.}
\begin{itemize}
  \item \textbf{Compressible and multiphysics.} Shocks, reacting flows, MHD ---
    regimes where rank behaviour is essentially unmapped.
  \item \textbf{The theory gap.} No sharp rank bounds for Navier--Stokes
    solutions as a function of $\mathrm{Re}$. Contrast this with
    \cref{subsec:mps}, where the one-dimensional area law~\cite{Hastings2007}
    gives exactly such a bound for quantum ground states. That is the missing
    foundation beneath every empirical result quoted above.
  \item \textbf{Hybrid tensor networks and machine learning.} Learned closures,
    learned preconditioners, tensor-network surrogates --- which is
    Part~\partref{4}.
  \item \textbf{Quantum hand-off.} The same MPS encoding is the natural
    interface to future quantum CFD algorithms.
\end{itemize}

\begin{exercise}[Companion notebooks]
  \texttt{03-walkthrough-forward-problem.ipynb} builds the solver above on a
  $64\times64$ grid; \texttt{03-exercises-forward-problem.ipynb} is the hands-on
  half-hour. Keep $n = 6$ unless a cell says otherwise, or you will spend the
  session watching a progress bar.

  \begin{enumerate}
    \item Instrument the advection term and watch the bond dimension of
      $u\,\partial_x u$ \emph{before} rounding. It is $\chi_u \chi_{\partial u}$,
      exactly. Then compare the advective and conservative forms: identical in
      exact arithmetic, not identical under truncation.
    \item Sweep $\chi_{\max}$ and find the smallest value that keeps
      $\norm{\nabla\cdot\mathbf{v}}$ below $10^{-1}$ for thirty steps --- then
      look at what that costs relative to the dense field. On a $64\times64$
      grid the answer is sobering, and the reason is that compression is an
      asymptotic statement.
    \item Put a smoothed cylinder in the flow. One extra Hadamard product per
      step, and no other change. Watch the sign convention on the level set: $q$
      must be \emph{negative} inside the solid, or you will faithfully simulate a
      flow that exists only inside the cylinder.
    \item Compress an intermittent field at several $\chi$ and compare the
      $\ell_2$ error, the spectrum and the flatness. Convince yourself, with your
      own numbers, which diagnostic you are allowed to trust.
  \end{enumerate}
\end{exercise}

\section{Part IV: The Inverse Problem - Data-Driven Approaches}
\label{sec:part4}
\lecture{01}{10}{2026}

% Mirrors slides.qmd "Part IV" outline. One subsection per outline item.

Parts~\partref{1}--\partref{3} assumed we know the governing equation
and want its solution. Part~\partref{4} inverts that assumption. What we
have instead are \textbf{snapshots} $\{u(t_k)\}$ --- from an experiment, from
DNS output, or from sensors. What we lack is the operator, the closure, and
sometimes the equation itself.

This is the regime of most real measurement data, and of every system too
complex to write down. The claim of this part is that the same compressed format
serves here too, and that it does so without giving up interpretability.

\subsection{Prediction of dynamics using time-series}
\label{subsec:prediction-time-series}
% Keywords: forecasting from snapshots, learning vs. prediction, snapshot
% matrices, curse of dimensionality in data-driven methods, why compressed
% formats help.

\paragraph{Two tasks, kept separate.}
\margintext{Conflating these two is the most common error in the data-driven
literature. A method can be excellent at one and useless at the other.}
\begin{itemize}
  \item \textbf{Prediction.} Given the past, forecast $u(t_{K+k})$. No
    interpretable model is required; only the forecast is judged.
  \item \textbf{Learning, or identification.} Recover the operator or the
    governing equations. Here interpretability is the entire point.
  \item Neural surrogates are the clean illustration: they predict well and
    explain nothing.
\end{itemize}

The classical toolkit spanning both tasks: POD and PCA, DMD and Koopman
operators, SINDy, operator learning (DeepONet, FNO), and data assimilation.

\paragraph{The same curse, again.}
\begin{itemize}
  \item Snapshot matrices are $N_x \times N_t$ with $N_x = 2^n$, so even taking
    the SVD is prohibitive --- the wall of
    \cref{subsec:numerical-methods-intro} in a new costume.
  \item The data brings its own problems on top: sparse and noisy measurements,
    short time series, non-stationarity, and no guarantee that a closed model
    exists at all.
  \item \textbf{The tensor-network promise:} compress \emph{both} space and time,
    perform the linear algebra inside the compressed format, and keep the modes
    interpretable.
\end{itemize}

\begin{remark}
  A useful practical bonus: a model obtained in MPS format plugs straight back
  into the Part~\partref{3} solver. Same data structure, no conversion
  layer, no lossy hand-off between the learned model and the simulator.
\end{remark}

\subsection{Dynamic Mode Decomposition and Space-Time Encoding}
\label{subsec:dmd-space-time}
% Keywords: DMD algorithm, Koopman operator viewpoint, snapshot matrices,
% space-time tensorization, QTT encoding of spatiotemporal data, connection
% to MPS-DMD.

\paragraph{DMD in one line.}
Assume the dynamics are approximately linear in the state,
\begin{equation}
  u_{k+1} \approx A\, u_k ,
  \label{eq:dmd-linear}
\end{equation}
find the best such $A$ from the data, and read off its eigenpairs. Each
eigenvalue gives a \textbf{growth rate and a frequency}; each eigenvector gives a
\textbf{spatial mode}. That is the whole appeal: the output is a set of
individually interpretable physical structures.

\paragraph{The standard recipe.}
\begin{itemize}
  \item Build snapshot matrices $X = [u_0 \cdots u_{T-1}]$ and
    $X' = [u_1 \cdots u_T]$.
  \item Take the SVD $X = U\Sigma V^{\dagger}$ and project onto the leading
    modes,
    \begin{equation}
      \tilde{A} = U^{\dagger} X' V \Sigma^{-1}.
      \label{eq:dmd-projected}
    \end{equation}
  \item Eigendecompose the small matrix $\tilde{A}$; the DMD modes are
    $\Phi = UW$.
\end{itemize}

\begin{remark}[Koopman view]
  Equation~\eqref{eq:dmd-projected} is a finite-rank approximation of the linear,
  infinite-dimensional Koopman operator acting on observables. That is why a
  linear method can describe nonlinear dynamics without contradiction.
\end{remark}

\paragraph{Space and time in a single MPS.}
\margintext{Our 2025 space--time paper. The conceptual move is small and the
consequence is large: once time is just more legs, there is no time loop to
parallelize or to accumulate error along.}
\begin{itemize}
  \item The idea~\cite{peddinti2025spacetime}: binarize \textbf{both} $x$ and
    $t$, into one MPS of $\log_2 N_x + \log_2 N_t$ tensors.
  \item Time stops being a loop and becomes simply more legs on the same chain.
  \item \textbf{Ordering matters here too.} Space-then-time works generally;
    time-then-space converged better for the nonlinear Schr\"odinger case,
    giving smoother vectorized solutions.
  \item The conceptual payoff: the \textbf{temporal} bond dimension becomes a
    direct measure of temporal complexity. Smooth dynamics means small $\chi$.
\end{itemize}

\paragraph{The all-at-once solve.}
\begin{itemize}
  \item Assemble the entire space--time linear system as a single MPO and solve
    it \textbf{once} with DMRG --- there is no time-stepping loop at all.
  \item Nonlinear terms are handled by Picard iteration, with linear
    convergence; Newton would give quadratic convergence at higher cost per
    iteration.
\end{itemize}

\noindent
Benchmarks on a $1024 \times 1024$ space--time grid:

\begin{center}
  \begin{tabular}{lll}
    \toprule
    \textsc{Equation} & $\chi$ & \textsc{Compression} \\
    \midrule
    Heat, wave                    & $2$--$22$ & $> 99\%$ \\
    Burgers, nonlinear diffusion  & $10$      & $> 99\%$, residual $\sim\!10^{-3}$ \\
    \bottomrule
  \end{tabular}
\end{center}

\paragraph{MPS-DMD.}
Standard DMD needs an SVD of an $N_x \times N_t$ matrix. That is the bottleneck
--- and, once the data is already a space--time MPS, it is unnecessary.

\begin{itemize}
  \item \textbf{MPS-DMD}~\cite{peddinti2025spacetime}: run the entire DMD
    pipeline \emph{inside} the space--time MPS, never forming the snapshot
    matrix.
  \item The key observation is that the MPS already contains the SVD. It only
    has to be gauged into the right form --- see \cref{subsec:mps}.
  \item Same modes, same eigenvalues, same predictions, at exponentially lower
    cost.
\end{itemize}

\noindent
The algorithm is three steps:

\begin{enumerate}
  \item \textbf{Gauge} into mixed canonical form at the space/time bond. The
    Schmidt decomposition there \emph{is} $U\Sigma V^{\dagger}$, so the SVD comes
    for free.
  \item \textbf{Project} the evolution operator onto the POD basis,
    $\tilde{A} = U^{\dagger} A U$: a small $\chi \times \chi$ matrix.
  \item \textbf{Eigendecompose} $\tilde{A} = W\Lambda W^{-1}$, take modes
    $\Phi = UW$, and predict by local contractions alone,
    \begin{equation}
      u_{T+k} = A^k u_T = U W \Lambda^k W^{-1} U^{\dagger} u_T .
      \label{eq:mps-dmd-predict}
    \end{equation}
\end{enumerate}

For externally supplied data --- where no space--time solve produced the MPS ---
stack the snapshot MPSs sequentially, truncating after each. The cost is linear
in the number of timesteps.

\paragraph{Results.}
\margintext{The headline is the scaling, not the compression ratio: runtime
logarithmic in \emph{both} $N_x$ and $N_t$, where prior tensor DMD variants are
linear in $N_t$.}
\begin{center}
  \begin{tabular}{@{}lp{0.62\textwidth}@{}}
    \toprule
    \textsc{Case} & \textsc{Setup and outcome} \\
    \midrule
    Burgers
      & $\chi = 10$ on $1024^2$; accurate $1000$ steps ahead, at $99\%$
        compression \\
    2D cylinder wake
      & $256\times128\times1024$ as a 25-bit MPS at $\chi = 200$; $97.42\%$
        compression, with 100 modes capturing the vortex street \\
    \bottomrule
  \end{tabular}
\end{center}

\subsection{Recovering Physics from Data}
\label{subsec:recovering-physics}
% Keywords: learning vs. prediction tasks, system identification, data
% assimilation, sparse/noisy measurements, governing-equation discovery.

DMD recovers an operator. The more ambitious task is to recover the
\emph{equation} --- and the difference matters, because an operator is tied to
the regime it was fitted in whereas an equation is not.

\begin{itemize}
  \item \textbf{System identification.} Fit a parametrized operator or a
    dictionary of candidate terms, and select among them. SINDy-style sparse
    regression is the canonical example; TT-based variants keep the dictionary
    manageable at high resolution.
  \item \textbf{Data assimilation.} Combine a known but imperfect model with
    sparse observations. Here the tensor network carries the state, and the
    model is the Part~\partref{3} solver.
  \item \textbf{Sparse and noisy measurements} are the normal case, not the
    exception. Recall from \cref{subsec:finding-qtt} that cross interpolation
    will fit noise faithfully if you let it.
  \item \textbf{Governing-equation discovery} is the hardest of these, and the
    one where interpretability is not a bonus but the deliverable.
\end{itemize}

\begin{remark}
  The truncation that makes these methods affordable is also, viewed
  differently, a regularizer: discarding small singular values discards
  low-amplitude structure, some of which is noise. Whether it denoises or
  faithfully reproduces the noise depends on where the noise sits in the
  spectrum, and that is not something you get to assume.
\end{remark}

\subsection{Tensor Networks for Learning Dynamical Laws}
\label{subsec:tn-learning-dynamics}
% Keywords: TN-based surrogate models, operator learning in TT format,
% training/optimization strategies, generalization and extrapolation in
% time, comparison to neural-network surrogates.

\begin{itemize}
  \item \textbf{TT surrogates and operator learning.} Represent the solution
    operator itself as an MPO and fit it to data. The training problem is again
    a sweeping optimization, so the machinery of
    \cref{subsec:linear-solvers} carries over directly.
  \item \textbf{Tensor-network closures.} Learn only the unresolved part --- the
    closure term --- and leave the resolved dynamics to the solver. This is the
    hybrid route, and the one where the physics prior does the most work.
  \item \textbf{Training strategies.} Sweeping optimization with adaptive ranks,
    warm starts across parameter values, and rank as the primary regularization
    knob.
  \item \textbf{Generalization and extrapolation.} Extrapolating in time is
    where surrogate models usually fail. A linear-in-observables model of the
    DMD kind extrapolates in a way you can reason about; a neural surrogate
    typically does not.
\end{itemize}

\begin{remark}[The comparison that matters]
  Against neural-network surrogates the tensor-network case does not rest on raw
  accuracy. It rests on the fact that modes, frequencies and growth rates come
  out of the method directly, rather than from a post-hoc probe --- and on the
  fact that rank gives a single, monitorable complexity knob.
\end{remark}

\subsection{Outlook}
\label{subsec:part4-outlook}
% Keywords: open challenges in data-driven TN methods, hybrid
% physics+data approaches, scalability, uncertainty quantification.

\paragraph{Extending the method.}
\margintext{The $(1+1)$D restriction is the most immediate limitation of the
space--time work, and lifting it brings back the index-ordering questions of
\cref{subsec:qtt-functions-operators}.}
\begin{itemize}
  \item \textbf{Beyond $(1+1)$D.} The space--time analysis so far is
    one-dimensional in space. Two and three dimensions plus time --- and the
    index ordering that goes with them --- is open.
  \item \textbf{Nonlinear and Koopman extensions.} Extended DMD with TT
    dictionaries, and kernel DMD in compressed form.
  \item \textbf{Noise and sparsity.} How TT truncation interacts with
    measurement noise is not settled; it can denoise, or it can fit the noise.
  \item \textbf{Online and streaming.} Update the MPS as data arrives, without
    recompressing from scratch.
\end{itemize}

\paragraph{Practice and positioning.}
\begin{itemize}
  \item \textbf{Hybrid physics and data.} Use the Part~\partref{3} solver as
    a prior and learn only the closure or the residual in TT form.
  \item \textbf{Uncertainty quantification.} Error bars on tensor-network
    predictions are essentially absent today. This is a serious gap for any
    engineering use.
  \item \textbf{Interpretability is the selling point.}
  \item \textbf{Software and reproducibility.} Shared TT/MPS tooling, standard
    benchmarks, and open datasets.
\end{itemize}



\printbibliography[heading=bibintoc] % heading=bibintoc adds the References section to the ToC
% shared bibliography.bib lives at the repo root, loaded via biblatex in preamble.tex

\clearpage
\appendix
\section{Appendix}
\label{sec:appendix}

% Mirrors slides.qmd "Appendix" outline. One subsection per outline item.

The two algorithms below are used throughout the notes. They are stated here in
full so the main text can refer to them without interrupting its argument.

\subsection{TT-Cross}
\label{subsec:appendix-tt-cross}
% Keywords: tensor cross interpolation algorithm details, maxvol/rectangular
% maxvol, skeleton decomposition, pivoting strategy, convergence and error
% bounds, computational cost.

TT-cross builds a tensor train for a black-box $f$ from $O(n\chi^2)$
evaluations, never forming the full array~\cite{oseledets2010ttcross}. It is
what makes \cref{subsec:finding-qtt,subsec:tt-cross-geometry} possible.

\paragraph{The matrix case first.}
\begin{itemize}
  \item A rank-$r$ matrix $M$ is exactly reconstructed from $r$ rows $I$ and $r$
    columns $J$ by the \textbf{skeleton} (or CUR) formula
    \begin{equation}
      M = M(:,J)\, M(I,J)^{-1} M(I,:) .
      \label{eq:skeleton}
    \end{equation}
  \item The whole question is which rows and columns to pick. A badly
    conditioned $M(I,J)$ destroys the reconstruction, since its inverse appears
    in~\eqref{eq:skeleton}.
  \item \textbf{maxvol} answers it: choose the $r \times r$ submatrix of maximum
    volume, i.e.\ largest $\abs{\det}$.\margintext{The maximum-volume principle
    is due to Goreinov and Tyrtyshnikov (2001); Oseledets and Tyrtyshnikov
    carried it up to tensor trains in 2010.} For an approximately rank-$r$ $M$
    the resulting error is within a factor $(r+1)$ of the best rank-$r$
    approximation in the Chebyshev norm.
  \item Finding the true maximum-volume submatrix is NP-hard. The practical
    algorithm swaps rows greedily until no single swap increases $\abs{\det}$,
    which lands within a controlled factor of the true maximum and costs
    $O(N r^2)$ per sweep.
\end{itemize}

\paragraph{The tensor-train algorithm.}
Maintain, for each bond $k$, a set of row multi-indices $\mathcal{I}_k$ and
column multi-indices $\mathcal{J}_k$. Then:

\begin{enumerate}
  \item Initialize $\mathcal{I}_k$ and $\mathcal{J}_k$, either at random or from
    a coarse pass over the domain.
  \item Sweeping left to right, form the slice
    $F_k = f(\mathcal{I}_k,\, i_k,\, \mathcal{J}_k)$, a $\chi \times d\chi$
    matrix obtained by varying only the local index $i_k$ while the pivots hold
    everything else fixed. This is the only place $f$ is called.
  \item Apply \textbf{maxvol} to $F_k$: select the $\chi$ rows of maximum
    $\abs{\det}$, and use them to update $\mathcal{I}_{k+1}$.
  \item Set the core to $F_k \hat{F}_k^{-1}$, where $\hat{F}_k$ is the selected
    maximum-volume submatrix, and carry the residual factor to the right.
  \item Sweep back right to left and repeat until the cross error, measured on
    fresh sample points, plateaus.
\end{enumerate}

\begin{remark}[Interpolation property]
  The resulting MPS is \textbf{exact} on every fibre it sampled --- that is what
  $\hat{F}_k^{-1}$ in step~4 enforces. This is the algorithm's defining
  strength and its defining weakness: it reproduces sampled values exactly,
  including sampled noise, and it says nothing at all about the fibres it never
  looked at.
\end{remark}

\paragraph{Practical notes.}
\begin{itemize}
  \item Cost is $O(n \chi^3)$ arithmetic and $O(n \chi^2)$ function evaluations
    for $n$ tensors of bond dimension $\chi$ --- \textbf{independent of $2^n$}.
  \item \textbf{Rectangular maxvol} keeps more than $\chi$ pivots, improving
    robustness at modest extra cost. Generally worth switching on.
  \item \textbf{Rank-adaptive variants} insert pivots where the local error is
    largest, rather than fixing $\chi$ in advance~\cite{nunez2025tci}.
  \item Convergence is heuristic. There is no guarantee for adversarial or noisy
    black boxes, and the failure modes of \cref{subsec:finding-qtt} are all
    observed in practice. Always validate on points the cross never sampled.
\end{itemize}

\subsection{DMRG}
\label{subsec:appendix-dmrg}
% Keywords: density matrix renormalization group, two-site sweep algorithm,
% local optimization/eigensolver, sweeping schedule, connection to
% variational MPS optimization, adaptive bond-dimension growth.

DMRG~\cite{white1992,schollwoeck2011} is the variational optimization of an MPS
by local updates. In these notes it appears as the linear solver of
\cref{subsec:linear-solvers}, as the pressure-Poisson step of
\cref{subsec:qtt-cfd-frameworks}, and as the all-at-once space--time solve of
\cref{subsec:dmd-space-time}.

\paragraph{The two-site sweep.}
\begin{enumerate}
  \item Bring the MPS into mixed canonical form with the orthogonality centre at
    site $k$ (\cref{subsec:mps}).
  \item Merge sites $k$ and $k+1$ into a single two-site tensor of size
    $\chi d^2 \chi$.
  \item Solve the \textbf{local problem} in that tensor alone, with every other
    tensor held fixed: an eigenproblem for ground states, a linear system for
    $Ax = b$. The effective operator is assembled from cached left and right
    \textbf{environments} --- the partial contractions of the MPO with the
    frozen parts of the chain --- which are updated one site at a time rather
    than rebuilt.
  \item Split the result by SVD, truncating to a tolerance $\varepsilon$ or a
    rank cap $\chi_{\max}$.
  \item Move the orthogonality centre one site along and repeat, sweeping right
    then left.
\end{enumerate}

\begin{remark}
  Because the canonical form makes the environments orthonormal, the local
  problem in step~3 is a genuine restriction of the global one, not an
  approximation of it. This is why the global objective decreases monotonically.
\end{remark}

\paragraph{Cost and practical notes.}
\begin{itemize}
  \item Per sweep, $O(n \chi^3 d^2 + n \chi^2 d\, \chi_{\mathrm{MPO}})$ for the
    contractions, with the local solve reaching $O(\chi^6)$ in the worst case.
    Observed totals are much better (\cref{subsec:linear-solvers}).
  \item \textbf{Adaptive bond dimension} is the point of the two-site variant:
    the SVD in step~4 may return a larger rank than it started with, so
    $\varepsilon$ rather than $\chi$ becomes the control parameter.
  \item \textbf{Escaping local minima.} Two-site updates explore directions
    unavailable to any single tensor, so they escape stalls that trap one-site
    sweeps. AMEn recovers most of this at one-site cost by enriching the update
    with residual directions~\cite{dolgov2014amen,holtz2012als}.
  \item \textbf{Sweeping schedule.} Start with a loose tolerance and a low rank
    cap, then tighten both over successive sweeps. Two or three sweeps usually
    suffice from a warm start.
  \item \textbf{Warm starts} matter more than any other tuning parameter in a
    time-stepping context, because consecutive solutions are close.
\end{itemize}

\end{document}
